% Options for packages loaded elsewhere
% Options for packages loaded elsewhere
\PassOptionsToPackage{unicode}{hyperref}
\PassOptionsToPackage{hyphens}{url}
\PassOptionsToPackage{dvipsnames,svgnames,x11names}{xcolor}
%
\documentclass[
  british,
  suppldata]{interact}
\usepackage{xcolor}
\usepackage{amsmath,amssymb}
\setcounter{secnumdepth}{-\maxdimen} % remove section numbering
\usepackage{iftex}
\ifPDFTeX
  \usepackage[T1]{fontenc}
  \usepackage[utf8]{inputenc}
  \usepackage{textcomp} % provide euro and other symbols
\else % if luatex or xetex
  \usepackage{unicode-math} % this also loads fontspec
  \defaultfontfeatures{Scale=MatchLowercase}
  \defaultfontfeatures[\rmfamily]{Ligatures=TeX,Scale=1}
\fi
\usepackage{lmodern}
\ifPDFTeX\else
  % xetex/luatex font selection
\fi
% Use upquote if available, for straight quotes in verbatim environments
\IfFileExists{upquote.sty}{\usepackage{upquote}}{}
\IfFileExists{microtype.sty}{% use microtype if available
  \usepackage[]{microtype}
  \UseMicrotypeSet[protrusion]{basicmath} % disable protrusion for tt fonts
}{}
\makeatletter
\@ifundefined{KOMAClassName}{% if non-KOMA class
  \IfFileExists{parskip.sty}{%
    \usepackage{parskip}
  }{% else
    \setlength{\parindent}{0pt}
    \setlength{\parskip}{6pt plus 2pt minus 1pt}}
}{% if KOMA class
  \KOMAoptions{parskip=half}}
\makeatother
% Make \paragraph and \subparagraph free-standing
\makeatletter
\ifx\paragraph\undefined\else
  \let\oldparagraph\paragraph
  \renewcommand{\paragraph}{
    \@ifstar
      \xxxParagraphStar
      \xxxParagraphNoStar
  }
  \newcommand{\xxxParagraphStar}[1]{\oldparagraph*{#1}\mbox{}}
  \newcommand{\xxxParagraphNoStar}[1]{\oldparagraph{#1}\mbox{}}
\fi
\ifx\subparagraph\undefined\else
  \let\oldsubparagraph\subparagraph
  \renewcommand{\subparagraph}{
    \@ifstar
      \xxxSubParagraphStar
      \xxxSubParagraphNoStar
  }
  \newcommand{\xxxSubParagraphStar}[1]{\oldsubparagraph*{#1}\mbox{}}
  \newcommand{\xxxSubParagraphNoStar}[1]{\oldsubparagraph{#1}\mbox{}}
\fi
\makeatother


\usepackage{longtable,booktabs,array}
\usepackage{calc} % for calculating minipage widths
% Correct order of tables after \paragraph or \subparagraph
\usepackage{etoolbox}
\makeatletter
\patchcmd\longtable{\par}{\if@noskipsec\mbox{}\fi\par}{}{}
\makeatother
% Allow footnotes in longtable head/foot
\IfFileExists{footnotehyper.sty}{\usepackage{footnotehyper}}{\usepackage{footnote}}
\makesavenoteenv{longtable}
\usepackage{graphicx}
\makeatletter
\newsavebox\pandoc@box
\newcommand*\pandocbounded[1]{% scales image to fit in text height/width
  \sbox\pandoc@box{#1}%
  \Gscale@div\@tempa{\textheight}{\dimexpr\ht\pandoc@box+\dp\pandoc@box\relax}%
  \Gscale@div\@tempb{\linewidth}{\wd\pandoc@box}%
  \ifdim\@tempb\p@<\@tempa\p@\let\@tempa\@tempb\fi% select the smaller of both
  \ifdim\@tempa\p@<\p@\scalebox{\@tempa}{\usebox\pandoc@box}%
  \else\usebox{\pandoc@box}%
  \fi%
}
% Set default figure placement to htbp
\def\fps@figure{htbp}
\makeatother



\ifLuaTeX
\usepackage[bidi=basic]{babel}
\else
\usepackage[bidi=default]{babel}
\fi
% get rid of language-specific shorthands (see #6817):
\let\LanguageShortHands\languageshorthands
\def\languageshorthands#1{}
\ifLuaTeX
  \usepackage[english]{selnolig} % disable illegal ligatures
\fi


\setlength{\emergencystretch}{3em} % prevent overfull lines

\providecommand{\tightlist}{%
  \setlength{\itemsep}{0pt}\setlength{\parskip}{0pt}}



 
\usepackage[]{natbib}
\bibliographystyle{tfcad}


\bibpunct[, ]{(}{)}{;}{a}{}{,}
\usepackage{booktabs}
\usepackage{longtable}
\usepackage{array}
\renewcommand{\thetable}{S\arabic{table}}
\makeatletter
\@ifpackageloaded{caption}{}{\usepackage{caption}}
\AtBeginDocument{%
\ifdefined\contentsname
  \renewcommand*\contentsname{Table of contents}
\else
  \newcommand\contentsname{Table of contents}
\fi
\ifdefined\listfigurename
  \renewcommand*\listfigurename{List of Figures}
\else
  \newcommand\listfigurename{List of Figures}
\fi
\ifdefined\listtablename
  \renewcommand*\listtablename{List of Tables}
\else
  \newcommand\listtablename{List of Tables}
\fi
\ifdefined\figurename
  \renewcommand*\figurename{Figure}
\else
  \newcommand\figurename{Figure}
\fi
\ifdefined\tablename
  \renewcommand*\tablename{Table}
\else
  \newcommand\tablename{Table}
\fi
}
\@ifpackageloaded{float}{}{\usepackage{float}}
\floatstyle{ruled}
\@ifundefined{c@chapter}{\newfloat{codelisting}{h}{lop}}{\newfloat{codelisting}{h}{lop}[chapter]}
\floatname{codelisting}{Listing}
\newcommand*\listoflistings{\listof{codelisting}{List of Listings}}
\makeatother
\makeatletter
\makeatother
\makeatletter
\@ifpackageloaded{caption}{}{\usepackage{caption}}
\@ifpackageloaded{subcaption}{}{\usepackage{subcaption}}
\makeatother
\usepackage{bookmark}
\IfFileExists{xurl.sty}{\usepackage{xurl}}{} % add URL line breaks if available
\urlstyle{same}
\hypersetup{
  pdflang={en-GB},
  colorlinks=true,
  linkcolor={blue},
  filecolor={Maroon},
  citecolor={Blue},
  urlcolor={Blue},
  pdfcreator={LaTeX via pandoc}}


\author{}
\date{}
\begin{document}

\articletype{SUPPLEMENTARY MATERIAL}
\title{Supplementary material for: Thermal time baselines for evaluating Sentinel-1, Sentinel-2 and PlanetScope predictors of wheat above-ground biomass under spatial cross-validation}
\author{\name{Francisco Zambrano\textsuperscript{a}, Ignacio Fuentes\textsuperscript{a}, Abel Herrera\textsuperscript{b} and Mauricio Molina-Rocco\textsuperscript{c}}
\affil{\textsuperscript{a}Facultad de Medicina Veterinaria y Agronom\'{i}a, Universidad de Las Am\'{e}ricas, Santiago, Chile; \textsuperscript{b}Hemera Centro de Observaci\'{o}n de la Tierra, Facultad de Ciencias, Ingenier\'{i}a y Tecnolog\'{i}a, Universidad Mayor, Santiago, Chile; \textsuperscript{c}Laboratorio de Suelos y Sostenibilidad, Departamento de Acuicultura y Recursos Agroalimentarios, Universidad de Los Lagos, Osorno, Chile}}
\maketitle


\section{Table S1: additional vegetation
indices}\label{table-s1-additional-vegetation-indices}

Table 3 of the main text lists the 21 vegetation indices (VIs), derived
from Sentinel-2 (S2) and PlanetScope (PS) spectral bands, that were
retained as predictors after eliminating highly correlated indices
(section 2.2.5 of the main text). Table S1 lists the remaining 32 VIs
that were part of the same 53-index base candidate set (section 2.2.5)
but were not retained in table 3.

As described in the main text, a subset of these S2-derived indices were
additionally recalculated using the Sentinel-2 Band 8A (narrow
near-infrared) instead of Band 8 (broad near-infrared), following the
same formula with the near-infrared (NIR) term replaced by the Band 8A
reflectance; this is denoted with an asterisk (*) next to the S2 dataset
code, following the same convention as table 3 of the main text.
Together, the 53 base indices (21 in table 3 and 32 in table S1) and
their 36 Band-8A-recalculated counterparts make up the full set of 89
candidate VIs referenced in section 2.2.5 of the main text.

Band and index notation follows table 3 of the main text: Red, Green,
Blue = visible bands; RE1, RE2, RE3 = S2 red-edge bands 5, 6, and 7; NIR
= near-infrared band (S2 Band 8 / PS Band 8); SWIR1, SWIR2 = S2
shortwave-infrared bands 11 and 12; Water vapour = S2 Band 9. Formula
definitions follow standard vegetation-index literature, including the
references cited in section 2.2.5 of the main text for VI selection
\citep{Peng2023, Segarra2022, Uribeetxebarria2023, Wang2022, Wu2008, Zheng2017};
the exact implementation of every index listed here and in table 3 is
available in the project code repository (script
\texttt{02\_calcular\_vis.R}).

\begingroup\scriptsize

\begin{longtable}{@{}c l >{\raggedright\arraybackslash}p{3.2cm} >{\raggedright\arraybackslash}p{2.0cm} >{\raggedright\arraybackslash}p{4.6cm}@{}}

\caption{\label{tbl-s1}Additional vegetation indices (covariates) derived from Sentinel-2 (S2) and PlanetScope (PS) spectral bands, computed as part of the 53-index base candidate set (section 2.2.5 of the main text) but not retained in table 3. Asterisk (*) next to the S2 dataset code indicates that the index was also recalculated using Sentinel-2 Band 8A (narrow NIR) instead of Band 8 (broad NIR).}

\tabularnewline

\\
\toprule
\textbf{N°} & \textbf{Dataset} & \textbf{Name} & \textbf{Abbreviation} & \textbf{Formula} \\
\midrule
\endfirsthead
\toprule
\textbf{N°} & \textbf{Dataset} & \textbf{Name} & \textbf{Abbreviation} & \textbf{Formula} \\
\midrule
\endhead
\bottomrule
\endlastfoot
1 & S2 & NDVI Variant (Red-edge 3) & NDVI$_{2}$ & $(RE_{3} - Red)/(RE_{3} + Red)$ \\
2 & S2 & Red-edge NDVI (705 nm) & NDVI$_{705}$ & $(RE_{2} - RE_{1})/(RE_{2} + RE_{1})$ \\
3 & S2*-PS & Modified Simple Ratio & MSR & $(NIR/Red - 1)/\sqrt{(NIR/Red + 1)}$ \\
4 & S2 & Red-edge Modified Simple Ratio & MSR$_{705}$ & $(RE_{2}/RE_{1} - 1)/\sqrt{(RE_{2}/RE_{1} + 1)}$ \\
5 & S2*-PS & Soil-Adjusted Vegetation Index & SAVI & $\frac{1.5 \cdot (NIR - Red)}{(NIR + Red + 0.5)}$ \\
6 & S2 & SAVI Variant (Red-edge 3) & SAVI$_{2}$ & $\frac{1.5 \cdot (RE_{3} - Red)}{(RE_{3} + Red + 0.5)}$ \\
7 & S2*-PS & Green SAVI & SAVI$_{green}$ & $\frac{1.5 \cdot (NIR - Green)}{(NIR + Green + 0.5)}$ \\
8 & S2*-PS & Green OSAVI & OSAVI$_{green}$ & $\frac{1.16 \cdot (NIR - Green)}{(NIR + Green + 0.16)}$ \\
9 & S2*-PS & Red-edge OSAVI & OSAVI$_{rededge}$ & $\frac{1.16 \cdot (NIR - RE_{1})}{(NIR + RE_{1} + 0.16)}$ \\
10 & S2 & - & TCARI/OSAVI$_{2}$ & $TCARI/{OSAVI}_{2}$ \\
11 & S2* & SWIR 11 Related OSAVI & OSAVI$_{SWIR11}$ & $\frac{1.16 \cdot (NIR - SWIR_{1})}{(NIR + SWIR_{1} + 0.16)}$ \\
12 & S2 & SWIR 12 Related TCARI & TCARI$_{SWIR12}$ & $3 \cdot ((RE_{1} - SWIR_{2}) - 0.2 \cdot (RE_{1} - Green) \cdot (RE_{1}/SWIR_{2}))$ \\
13 & S2* & SWIR 12 Related OSAVI & OSAVI$_{SWIR12}$ & $\frac{1.16 \cdot (NIR - SWIR_{2})}{(NIR + SWIR_{2} + 0.16)}$ \\
14 & S2*-PS & Chlorophyll Index - Green & CI$_{green}$ & $(NIR/Green) - 1$ \\
15 & S2*-PS & Chlorophyll Index - Red-edge & CI$_{rededge}$ & $(NIR/RE_{1}) - 1$ \\
16 & S2*-PS & OSAVI $\times$ Chlorophyll Index Red-edge & OSAVI $\times$ CI$_{rededge}$ & $\left[\frac{1.16 \cdot (NIR - Red)}{(NIR + Red + 0.16)}\right] \cdot \left[(NIR/RE_{1}) - 1\right]$ \\
17 & S2*-PS & Difference Vegetation Index & DVI & $NIR - Red$ \\
18 & S2*-PS & Green DVI & DVI$_{green}$ & $NIR - Green$ \\
19 & S2*-PS & Red-edge DVI & DVI$_{rededge}$ & $NIR - RE_{1}$ \\
20 & S2 & Two-Band Enhanced Vegetation Index & EVI2 & $\frac{2.5 \cdot (RE_{3} - Red)}{(1 + RE_{3} + 2.4 \cdot Red)}$ \\
21 & S2*-PS & Red-edge Vegetation Index 1 & REVI1 & $NIR - RE_{1}$ \\
22 & S2* & Red-edge Vegetation Index 2 & REVI2 & $NIR - RE_{2}$ \\
23 & S2*-PS & Weighted Difference Vegetation Index & WDVI & $NIR - 0.5 \cdot Red$ \\
24 & S2-PS & Green-Red Vegetation Index & GRVI & $(Green - Red)/(Green + Red)$ \\
25 & S2*-PS & Green Ratio Vegetation Index & GRVI2 & $NIR/Green$ \\
26 & S2*-PS & Green Atmospherically Resistant Index & GARI & $\frac{NIR - (Green - 1.7 \cdot (Blue - Red))}{NIR + (Green - 1.7 \cdot (Blue - Red))}$ \\
27 & S2*-PS & Leaf Chlorophyll Index & LCI & $(NIR - RE_{1})/(NIR - Red)$ \\
28 & S2 & MERIS Terrestrial Chlorophyll Index & MTCI & $(RE_{2} - RE_{1})/(RE_{1} - Red)$ \\
29 & S2* & Normalized Difference Red-edge Index 2 & NDRE2 & $(NIR - RE_{2})/(NIR + RE_{2})$ \\
30 & S2 & Normalized Difference Red-edge Index (RE3/Green) & NDREI & $(RE_{3} - Green)/(RE_{3} + Green)$ \\
31 & S2* & Normalized Difference Moisture Index & NDMI & $(NIR - SWIR_{1})/(NIR + SWIR_{1})$ \\
32 & S2*-PS & Kernel NDVI & KNDVI & $\frac{1 - \exp\left(-\frac{(NIR-Red)^2}{2 \cdot 0.15^2}\right)}{1 + \exp\left(-\frac{(NIR-Red)^2}{2 \cdot 0.15^2}\right)}$ \\

\end{longtable}

\endgroup

\section{Equations S1 to S3: performance
metrics}\label{equations-s1-to-s3-performance-metrics}

Section 2.3.4 of the main text reports model performance using the mean
absolute error (MAE), root mean square error (RMSE), and coefficient of
determination (R²), defined in equations (S1) to (S3) below and computed
separately for each held-out site-season, where \(y_i\) and
\(\hat{y}_i\) are the observed and out-of-fold predicted AGB values for
sample \(i\), \(\bar{y}\) is the mean observed AGB of the held-out
site-season, and \(n\) is the number of samples in that site-season. R²
computed in this way is negative when the predictions are less accurate
than the mean of the held-out site-season.

\[
MAE = \frac{1}{n}\sum_{i=1}^{n}\left|y_i - \hat{y}_i\right| \tag{S1}
\]

\[
RMSE = \sqrt{\frac{1}{n}\sum_{i=1}^{n}\left(y_i - \hat{y}_i\right)^2} \tag{S2}
\]

\[
R^2 = 1 - \frac{\displaystyle\sum_{i=1}^{n}\left(y_i - \hat{y}_i\right)^2}{\displaystyle\sum_{i=1}^{n}\left(y_i - \bar{y}\right)^2} \tag{S3}
\]

\section{Table S2: performance per held-out
site-season}\label{table-s2-performance-per-held-out-site-season}

Table S2 reports the RMSE and R² of every predictor set and model for
each of the four held-out site-seasons. Mean and standard deviation
across site-seasons are given in table 5 of the main text. The values
are produced by \texttt{script/31\_loso\_diseno\_limpio.R} and stored in
\texttt{output/tables/rev2\_loso\_limpio\_por\_fold.csv}.

\begin{table}[htbp]
\centering
\footnotesize
\caption{RMSE (t ha$^{-1}$) and, in parentheses, R$^2$ (1 $-$ SSE/SST) for each held-out site-season under leave-one-site-season-out cross-validation (LOSSO-CV). Each value is obtained with a model trained on the other three site-seasons.}
\label{tbl-s2}
\begin{tabular}{l l c c c c}
\toprule
\textbf{Predictor set} & \textbf{Model} & \textbf{Hidango North} & \textbf{Hidango South} & \textbf{La Cancha} & \textbf{Villa Baviera} \\
\midrule
DAS & LM & 5.57 (0.74) & 8.81 ($-$0.66) & 9.95 (0.44) & 7.04 (0.25) \\
GDD & LM & 5.80 (0.71) & 6.85 ($-$0.01) & 8.87 (0.56) & 3.51 (0.81) \\
GDD & RF & 8.48 (0.39) & 10.72 ($-$1.46) & 10.67 (0.36) & 5.19 (0.59) \\
GDD & XGBoost & 8.56 (0.38) & 10.87 ($-$1.53) & 10.73 (0.35) & 5.22 (0.59) \\
W & LM & 6.72 (0.61) & 7.17 ($-$0.10) & 8.97 (0.55) & 3.62 (0.80) \\
W & RF & 8.23 (0.42) & 7.69 ($-$0.27) & 7.86 (0.65) & 6.36 (0.39) \\
W & XGBoost & 7.51 (0.52) & 8.84 ($-$0.68) & 8.23 (0.62) & 6.46 (0.37) \\
S1 & RF & 12.48 ($-$0.33) & 10.56 ($-$1.39) & 12.67 (0.10) & 7.99 (0.04) \\
S1 & XGBoost & 12.74 ($-$0.38) & 11.24 ($-$1.71) & 14.49 ($-$0.18) & 10.29 ($-$0.60) \\
S2 & RF & 9.01 (0.31) & 5.57 (0.34) & 9.82 (0.46) & 3.34 (0.83) \\
S2 & XGBoost & 9.76 (0.19) & 5.68 (0.31) & 9.42 (0.50) & 4.04 (0.75) \\
PS & RF & 6.61 (0.63) & 4.91 (0.48) & 8.45 (0.60) & 14.62 ($-$2.23) \\
PS & XGBoost & 7.68 (0.50) & 5.36 (0.39) & 9.26 (0.52) & 15.57 ($-$2.66) \\
S1+GDD & RF & 5.43 (0.75) & 8.05 ($-$0.39) & 7.91 (0.65) & 3.87 (0.77) \\
S1+GDD & XGBoost & 5.65 (0.73) & 8.58 ($-$0.58) & 8.73 (0.57) & 3.88 (0.77) \\
S2+GDD & RF & 8.20 (0.43) & 5.43 (0.37) & 9.03 (0.54) & 2.39 (0.91) \\
S2+GDD & XGBoost & 9.30 (0.26) & 5.34 (0.39) & 8.47 (0.60) & 2.91 (0.87) \\
PS+GDD & RF & 5.61 (0.73) & 4.62 (0.54) & 7.99 (0.64) & 11.57 ($-$1.02) \\
PS+GDD & XGBoost & 6.68 (0.62) & 4.53 (0.56) & 8.58 (0.59) & 12.40 ($-$1.32) \\
S1+W & RF & 6.87 (0.60) & 7.51 ($-$0.21) & 6.20 (0.78) & 4.45 (0.70) \\
S1+W & XGBoost & 5.97 (0.70) & 8.27 ($-$0.46) & 7.38 (0.69) & 4.61 (0.68) \\
S2+W & RF & 8.42 (0.39) & 5.55 (0.34) & 8.35 (0.61) & 2.92 (0.87) \\
S2+W & XGBoost & 8.69 (0.36) & 5.69 (0.31) & 7.72 (0.66) & 2.85 (0.88) \\
PS+W & RF & 6.68 (0.62) & 5.05 (0.45) & 7.36 (0.70) & 10.20 ($-$0.57) \\
PS+W & XGBoost & 7.26 (0.55) & 5.00 (0.46) & 7.78 (0.66) & 7.92 (0.05) \\
S2+W+S1 & RF & 8.25 (0.42) & 5.54 (0.34) & 8.25 (0.62) & 2.23 (0.92) \\
S2+W+S1 & XGBoost & 8.66 (0.36) & 5.75 (0.29) & 8.04 (0.64) & 2.82 (0.88) \\
All & RF & 7.59 (0.51) & 5.02 (0.46) & 7.98 (0.64) & 3.85 (0.78) \\
All & XGBoost & 7.92 (0.47) & 5.38 (0.38) & 7.66 (0.67) & 3.66 (0.80) \\
\bottomrule
\end{tabular}
\end{table}


\section{Note S1: predictors removed by the correlation filter in each
fold}\label{note-s1-predictors-removed-by-the-correlation-filter-in-each-fold}

The correlation filter (Pearson r \textgreater{} 0.9) was estimated on
the training data of each LOSSO-CV fold, after imputation and
normalisation, so the retained predictors differ among folds. Table S3
lists the predictors removed from the set with all 54 candidate
predictors (All) in each fold. Most removed predictors are daily
Sentinel-2 indices that are highly correlated with one another.
\(\sum\)GDD was excluded from the filter, because it is highly
correlated with cumulative vegetation indices and would otherwise have
been removed in some folds. The removals can be reproduced with
\texttt{script/31\_loso\_diseno\_limpio.R}.

\begin{table}

\caption{\label{tbl-s3}Predictors removed by the within-fold correlation filter from the All set, by held-out field. The suffix 8A marks Sentinel-2 indices computed with band B8A, and cumsum marks cumulative sums from the start of the season.}

\centering{

\centering
\footnotesize


\begin{tabular}{l c p{0.62\linewidth}}
\toprule
\textbf{Held-out field} & \textbf{n} & \textbf{Removed predictors} \\
\midrule
Hidango North & 19 & S2: MCARI/OSAVI, MCARI/OSAVI2, TCARI/OSAVI2, OSAVI, OSAVI2, NDRE1, EVI, GNDVI, MCARI/OSAVI 8A, CI$_{red}$ 8A, SIPI 8A, MCARI$_{SW11}$, NDRE2, SR, TCARI/OSAVI, B1 cumsum, CI$_{red}$ 8A cumsum; PS: OSAVI, CI$_{red}$ \\
Hidango South & 16 & S2: TCARI/OSAVI2, OSAVI, OSAVI2, NDRE1, NDRE/NDVI, EVI, GNDVI, MCARI/OSAVI 8A, NDRE3 8A cumsum, MCARI/OSAVI, SR, TCARI/OSAVI, CI$_{red}$, SIPI; PS: OSAVI, CI$_{red}$ \\
La Cancha & 18 & S2: TCARI/OSAVI2, OSAVI, OSAVI2, NDRE1, EVI, GNDVI, CI$_{red}$, MCARI/OSAVI 8A, NDRE3 8A cumsum, TCARI, MCARI/OSAVI, MCARI$_{SW11}$, TCARI/OSAVI, MCARI, SR, SIPI; PS: OSAVI, SR \\
Villa Baviera & 24 & S2: MCARI/OSAVI, MCARI/OSAVI2, TCARI/OSAVI2, OSAVI, OSAVI2, NDRE1, EVI, GNDVI, MCARI/OSAVI 8A, CI$_{red}$ 8A, SR, TCARI/OSAVI, MCARI, SIPI, B1 cumsum, CI$_{red}$ 8A cumsum; PS: B8, SR, OSAVI, EVI, SIPI, CI$_{red}$, EVI cumsum, CI$_{red}$ cumsum \\
\bottomrule
\end{tabular}

}

\end{table}%


\bibliography{references.bib}



\end{document}
